\documentclass[12pt]{article}
\usepackage[utf8]{inputenc}
\usepackage[T1]{fontenc}
\usepackage{hyperref}
\title{Wardandi Language Resources for K-12 Curriculum Development}
\author{Compiled by Claude Sonnet 5.5 (claude-sonnet-5-5) for Andr\'as Kornai}
\date{}
\begin{document}
\maketitle

\section{Introduction}
Wardandi (ISO 639-3 wxw) is the language of the Wardandi people of south-western Western Australia. Their country runs from Capel in the north-east to Augusta in the south and includes Busselton, Yallingup and Margaret River. It is listed as extinct. The Noongar Boodjar Language Centre places Wardandi within the south-western ``Kongal-marawar'' dialect of Noongar. I found no Wardandi-specific dictionary, course or school programme, so the Noongar revival sources below are regional context.

\section{Primary Educational Resources}
The Western Australian curriculum includes a Noongar language revival course for pre-primary to Year 10. Two of its pages, Years 5 and 7, appeared in the searches. \url{https://k10outline.scsa.wa.edu.au/home/teaching/curriculum-browser/languages/western-australian-aboriginal-languages-noongar-language-revival-pre-primary-to-year-10/year-5/understanding/role-of-language-building/aclfwu172} and \url{https://k10outline.scsa.wa.edu.au/home/teaching/curriculum-browser/languages/western-australian-aboriginal-languages-noongar-language-revival-pre-primary-to-year-10/year-7/understanding/role-of-language-building/aclfwu194} These are Noongar materials, not Wardandi-specific ones.

\section{Linguistic and Reference Materials}
\begin{itemize}
\item AIATSIS AUSTLANG entry w3. \url{https://collection.aiatsis.gov.au/austlang/language/w3}
\item Bracknell et al., a paper from the AIATSIS Yumi Sabe project. \url{https://yumi-sabe.aiatsis.gov.au/sites/default/files/outputs/2024-02/Bracknell_etal.pdf}
\item Rebuilding as Research: Noongar Song, Language and Ways of Knowing (University of Western Australia). \url{https://research-repository.uwa.edu.au/en/publications/rebuilding-as-research-noongar-song-language-and-ways-of-knowing/}
\end{itemize}

\section{Religious and Community Resources}
\begin{itemize}
\item The Noongar Boodjar Language Centre in Bunbury works on recording, transcribing and archiving Noongar. I saw no direct URL for it, so it is not in the JSON.
\item Wardandi Clan Plaque page on the Noongar Boodjar site. \url{https://noongarboodjar.com.au/?p=1964}
\item ``My journey on coastal Noongar Boodjar'' (IYIL 2019). \url{https://en.iyil2019.org/wp-content/uploads/2019/05/my-journey-on-coastal-noongar-boodjar-noongar-country.pdf}
\item The Wirlomin Noongar group has published bilingual books for the south coast. \url{https://ro.ecu.edu.au/exhibit/highlighting-aboriginal-and-torres-strait-islander-research-at-ecu/language/ever-widening-circles-consolidating-and-enhancing-wirlomin-noongar-archival-material-in-the-community/}
\end{itemize}

\section{Digital and Audio Resources}
No Wardandi audio, video or software was found in the searches.

\section{Curriculum Development Context}
Wardandi has no fluent speakers, so any curriculum would be heritage and revitalization work. It would need to be developed with Wardandi and Noongar custodians. The WA Noongar language revival course is a framework that could be adapted, and the Noongar Boodjar Language Centre is the obvious organisation to ask.
\end{document}